\chapter{The one-dimensional area law and matrix product state approximation}
\label{ch:one_dimensional_chain}

% The chain specialization of the 2026 area-law and PEPS manuscripts, following the
% TNLean note docs/notes/one_dimensional_note. Every constant is explicit. The model is
% stated natively on the chain of L sites; nothing is derived from the two-dimensional
% lattice development.

The chain of $L$ sites carries a local dimension $q$, a finite-range Hamiltonian
with one Hermitian term of norm at most $J$ per support of diameter at most $R$,
and a unit ground vector separated from the rest of the spectrum by a gap
$\Delta$. On a chain the two-dimensional argument of \cite{OpenAI2026AreaLaw}
collapses to three steps: the marginal tail bound at an interval cut, the
buffered discount with a buffer of constant width, and the Araki--Lieb
inequality. The matrix product state approximation of
\cite{OpenAI2026PolynomialPEPS} then follows from the tail bound alone by
truncating the Schmidt decompositions across the cuts in sequence. This chapter
follows \cite{TNLean2026OneDimensional}; all constants are explicit, and the
area-law constant is linear in $1+J/\Delta$. The two finite-dimensional cores,
the tail bound and the buffered discount, are stated in the companion library
for an arbitrary finite set of sites and are instantiated here directly.

\section{Chains, supports and Hamiltonians}

\begin{definition}[Chains, intervals and finite-range Hamiltonians]
    \label{def:chain_model}
    \lean{
        TNLean.MPS.AreaLaw.interval,
        TNLean.MPS.AreaLaw.mem_interval,
        TNLean.MPS.AreaLaw.interval_mono,
        TNLean.MPS.AreaLaw.interval_eq_empty,
        TNLean.MPS.AreaLaw.interval_zero_eq_univ,
        TNLean.MPS.AreaLaw.card_interval_le,
        TNLean.MPS.AreaLaw.IsChainSupport,
        TNLean.MPS.AreaLaw.ChainSupport,
        TNLean.MPS.AreaLaw.IsChainSupport.card_le,
        TNLean.MPS.AreaLaw.supportDim_le,
        TNLean.MPS.AreaLaw.one_le_supportDim,
        TNLean.MPS.AreaLaw.ChainHamiltonian,
        TNLean.MPS.AreaLaw.ChainHamiltonian.operator,
        TNLean.MPS.AreaLaw.ChainHamiltonian.operator_isHermitian,
        TNLean.MPS.AreaLaw.ChainHamiltonian.zero,
        TNLean.MPS.AreaLaw.cutLogBudget,
        TNLean.MPS.AreaLaw.one_le_cutLogBudget
    }
    \leanok
    The \emph{chain of $L$ sites} has sites $0,\ldots,L-1$, each carrying
    $\mathbb C^q$. For natural numbers $a\le b$ the \emph{interval} $[a,b)$ is
    the set of sites with position at least $a$ and less than $b$; it has at
    most $b-a$ sites, and its dilation $[a-d,b+d)$ is clipped at the ends of the
    chain. A \emph{support of range $R$} is a nonempty set of sites of diameter
    at most $R$; it has at most $R+1$ sites and Hilbert-space dimension at most
    $q^{R+1}$. A \emph{finite-range Hamiltonian} assigns to every support $X$
    of range $R$ a Hermitian term $h_X$ acting as the identity off $X$ with
    operator norm at most $J$, zero terms allowed, and
    $H=\sum_X h_X$. A \emph{gapped ground vector} with gap $\Delta$ is a unit
    vector $\Omega$ with $H\Omega=E_0\Omega$ and
    $H-E_0\ge\Delta(1-|\Omega\rangle\langle\Omega|)$. For a cut $A$ the
    \emph{cut budget} is
    $\mathcal B_A=1+\sum_X\log^2(e\,q^{|X|})$ over the supports $X$ meeting
    both $A$ and its complement; it is at least one.
\end{definition}

\begin{lemma}[Marginal tail bound at a cut]
    \label{lem:chain_marginal_tail}
    \lean{TNLean.MPS.AreaLaw.ChainHamiltonian.surprisalTail_regionState_le}
    \leanok
    \uses{def:chain_model}
    Let $\Omega$ be a gapped ground vector, with gap $\Delta>0$, of a
    finite-range Hamiltonian with term norms at most $J\ge0$ on the chain, and
    let $A$ be any set of sites. Let $K$ be the surprisal $-\log p$ of the
    eigenvalue $p$ of the reduced state on $A$ drawn with probability $p$, so
    that $\mathbb E K=S(A)$. For every real $w$,
    \begin{align}
        \Pr\{|K-S(A)|>w\}\le
        \min\Bigl\{1,\ 2e^{e/2}\exp\Bigl(-\frac{w}{32\sqrt{(1+J/\Delta)\,
                \mathcal B_A}}\Bigr)\Bigr\}.
    \end{align}
\end{lemma}
\begin{proof}
    \leanok
    \uses{def:chain_model}
    This is Lemma~3.1 of \cite{OpenAI2026AreaLaw}, which holds on any finite
    tensor product: each term acts as the identity off its support, is
    Hermitian and has norm at most $J$, and the ground vector is a unit
    eigenvector with the projector gap. The chain model supplies exactly these
    hypotheses.
\end{proof}

\begin{lemma}[Interactions split by an interval]
    \label{lem:chain_crossing_count}
    \lean{
        TNLean.MPS.AreaLaw.IsChainSupport.subset_window_of_splits,
        TNLean.MPS.AreaLaw.chainCrossConst,
        TNLean.MPS.AreaLaw.card_filter_subset_le,
        TNLean.MPS.AreaLaw.card_crossingTerms_interval_le,
        TNLean.MPS.AreaLaw.chainBudgetConst,
        TNLean.MPS.AreaLaw.one_le_chainBudgetConst,
        TNLean.MPS.AreaLaw.supportDim_le_real,
        TNLean.MPS.AreaLaw.cutLogBudget_interval_le,
        TNLean.MPS.AreaLaw.not_splits_two_shells
    }
    \leanok
    \uses{def:chain_model}
    Let $X=[a,b)$ be an interval of the chain. A support of range $R$ that
    meets both $X$ and its complement lies in the window $[a-R,a+R)$ or in the
    window $[b-R,b+R)$, each of at most $2R$ sites. Hence at most
    \begin{align}
        N_R=2\cdot2^{2R}
    \end{align}
    supports are split by $X$, and the cut budget of $X$ is at most
    \begin{align}
        B_R=1+N_R\log^2\bigl(e\,q^{R+1}\bigr).
    \end{align}
    Both bounds are independent of the interval, and no support of range $R$
    is split by two of the shells $[l-(j+1)R,\,r+(j+1)R)$, $j\ge0$, of an
    interval $[l,r)$.
\end{lemma}
\begin{proof}
    \leanok
    \uses{def:chain_model}
    A split support contains a site $v$ in $X$ and a site $w$ outside $X$
    with $|v-w|\le R$. If $w<a$ then every site $x$ of the support satisfies
    $a-R\le v-R\le x\le w+R<a+R$; if $w\ge b$ then $b-R\le w-R\le x\le v+R<b+R$.
    A window of at most $2R$ sites has at most $2^{2R}$ subsets. The budget
    bound follows from the crossing count and the dimension bound $q^{R+1}$. If
    a support met the shell with index $j$ and left the shell with index
    $j'>j$, two of its sites would be at distance at least $R+1$.
\end{proof}

\begin{theorem}[Logarithmic dependence of the interval cut budget]
    \label{thm:chain_budget_log_sq}
    \lean{TNLean.MPS.AreaLaw.chainBudgetConst_le_mul_log_sq}
    \leanok
    \uses{lem:chain_crossing_count}
    Let $R$ be a nonnegative integer and put $N_R=2\cdot2^{2R}$ and
    $B_R(q)=1+N_R\log^2(e q^{R+1})$. For every integer $q\ge2$,
    \begin{align}
        B_R(q)\le\frac{B_R(2)}{(\log2)^2}(\log q)^2.
    \end{align}
    Here $N_R$ is the interval crossing bound for the
    family of all admissible supports, as in the cut-budget definition of
    \cite[Part~I]{TNLean2026OneDimensional}.
    The remark on sharper constants in that part describes its application
    to the logarithmic dependence of the area-law constant.
\end{theorem}
\begin{proof}
    \leanok
    \uses{lem:chain_crossing_count}
    Put $c=\log2>0$ and $t=\log q\ge c$. The identity
    $\log(e q^{R+1})=1+(R+1)t$ gives
    $c(1+(R+1)t)\le t(1+(R+1)c)$.
    Both sides are nonnegative. Squaring, multiplying by $N_R$ and adding
    $c^2\le t^2$ yields
    \begin{align}
        c^2 B_R(q)
         & =c^2+N_R\bigl(c(1+(R+1)t)\bigr)^2 \notag \\
         & \le t^2+N_R\bigl(t(1+(R+1)c)\bigr)^2
        =t^2 B_R(2). \notag
    \end{align}
    Divide by $c^2>0$.
\end{proof}


\section{The buffered discount with a constant-width buffer}

\begin{lemma}[Buffered discount on a chain]
    \label{lem:chain_buffer_discount}
    \lean{
        TNLean.MPS.AreaLaw.chainEnergyThreshold,
        TNLean.MPS.AreaLaw.chainTailThreshold,
        TNLean.MPS.AreaLaw.chainShellCount,
        TNLean.MPS.AreaLaw.four_le_chainShellCount,
        TNLean.MPS.AreaLaw.chainEnergyThreshold_le_chainShellCount,
        TNLean.MPS.AreaLaw.chainTailThreshold_le_chainShellCount,
        TNLean.MPS.AreaLaw.chainDiscountConst,
        TNLean.MPS.AreaLaw.chainBuffer,
        TNLean.MPS.AreaLaw.disjoint_interval_chainBuffer,
        TNLean.MPS.AreaLaw.card_chainBuffer_le,
        TNLean.MPS.AreaLaw.two_mul_condEntropy_interval_le
    }
    \leanok
    \uses{def:chain_model}
    Let $\Omega$ be a gapped ground vector, with gap $\Delta>0$, of a
    finite-range Hamiltonian with term norms at most $J\ge0$ on the chain with
    local dimension $q\ge1$. Put $\vartheta=J/\Delta$ and
    \begin{align}
        m=\Bigl\lceil\max\bigl\{4,\ 32N_Rq^{2(R+1)}\vartheta,\
        128\sqrt{(1+\vartheta)B_R}\bigr\}\Bigr\rceil .
    \end{align}
    For an interval $X=[l,r)$ let $T=X^{+mR}\setminus X$ be the buffer,
    which has at most $2mR$ sites. Then
    \begin{align}
        2\bigl(S(X\cup T)-S(T)\bigr)\le S(X)+
        \frac{4096\,e\,\vartheta B_R}{m}+m\log2 .
    \end{align}
    The energy threshold $\Delta/2$, the overlap losses $m\log2$ and the
    explicit tail term, with the three-entry $m$, replace the printed
    constants of Lemma~4.2 of \cite{TNLean2026OneDimensional}, which runs at
    the threshold $\Delta/4$, absorbs the tail term through a fourth entry of
    $m$ and concludes with $1+16N_Rq^{2(R+1)}\vartheta$; see
    \cite{gap:tnlean_one_dimensional_discount_constants}.
\end{lemma}
\begin{proof}
    \leanok
    \uses{lem:chain_crossing_count, def:chain_model}
    This is the finite-dimensional buffered discount (optimized filters,
    stationarity, the energy estimate, the prefix comparison and Schmidt
    pinning) of Lemma~3.2 of \cite{OpenAI2026AreaLaw}, applied to the nested
    shells $X_j=X^{+(j+1)R}$, $0\le j<m$, with the uniform weights $a_j=2/m$ of
    total mass two. The shells are intervals, so each is split by at most $N_R$
    supports with cut budget at most $B_R$ by
    Lemma~\ref{lem:chain_crossing_count}, and no support is split by two of
    them. The choice of $m$ makes the energy of the filtered vector at most
    $\Delta/2$, each weight at most $1/2$, and each tilt $a_j/(1-a_j)$ at most
    the tail radius $1/(32\sqrt{(1+\vartheta)B_j})$. The error terms sum to
    $1024e\vartheta\sum_jB_ja_j^2\le4096e\vartheta B_R/m$, and the overlap
    losses to $m\log2$.
\end{proof}

\section{Closing the argument}

\begin{lemma}[Araki--Lieb and the area law from interval bounds]
    \label{lem:chain_entropy_inequalities}
    \lean{
        TNLean.MPS.AreaLaw.regionEntropy_union_le,
        TNLean.MPS.AreaLaw.regionEntropy_le_union_add,
        TNLean.MPS.AreaLaw.regionEntropy_empty,
        TNLean.MPS.AreaLaw.regionEntropy_univ,
        TNLean.MPS.AreaLaw.regionEntropy_le_card_mul_log,
        TNLean.MPS.AreaLaw.boundaryEdges,
        TNLean.MPS.AreaLaw.mem_boundaryEdges,
        TNLean.MPS.AreaLaw.boundaryEdges_empty,
        TNLean.MPS.AreaLaw.boundaryEdges_univ,
        TNLean.MPS.AreaLaw.exists_interval_subset_boundaryEdges_sdiff_ssubset,
        TNLean.MPS.AreaLaw.regionEntropy_le_mul_card_boundaryEdges_of_intervals
    }
    \leanok
    \uses{def:chain_model}
    For a unit vector $\Omega$ on the chain and disjoint sets of sites $X,T$,
    $S(X\cup T)\le S(X)+S(T)$ and $S(X)\le S(X\cup T)+S(T)$; the empty set and
    the whole chain have entropy zero, and $S(A)\le|A|\log q$. The
    \emph{boundary edges} of a set $A$ are the edges $\{p,p+1\}$ of the chain
    with exactly one endpoint in $A$. If every interval has entropy at most
    $C\ge0$, then every subset $A$ has entropy at most $C\,|\partial A|$.
\end{lemma}
\begin{proof}
    \leanok
    \uses{def:chain_model}
    Subadditivity and purity are the finite-product entropy inequalities of the
    companion library; the Araki--Lieb inequality is purity,
    $S(X)=S(\Lambda\setminus X)$, followed by subadditivity for the disjoint
    union $\Lambda\setminus X=(\Lambda\setminus(X\cup T))\cup T$. The dimension
    bound is the entropy bound $\log\dim$ for a trace-one positive state. For
    the last claim, induct on the number of sites of $A$. If $A$ is empty or
    the whole chain its entropy vanishes. Otherwise split off the maximal run
    $I$ of consecutive sites starting at the least position of $A$.
    Subadditivity gives $S(A)\le S(I)+S(A\setminus I)\le C+S(A\setminus I)$,
    every boundary edge of $A\setminus I$ is a boundary edge of $A$, and $A$
    has at least one more boundary edge (the edge at the right end of $I$ when
    $A\setminus I$ is nonempty, and an edge at an end of $I$ otherwise, since
    $I$ is then a proper nonempty subset of the chain).
\end{proof}

\begin{theorem}[Entropy at local dimension one]
    \label{thm:chain_entropy_dim_one}
    \lean{TNLean.MPS.AreaLaw.regionEntropy_eq_zero_of_localDimension_one}
    \leanok
    \uses{def:chain_model}
    Let $\Omega$ be a unit vector on a finite chain with one-dimensional
    local spaces. For every set $A$ of sites, $S_\Omega(A)=0$. This is the
    dimension-one case of the entropy dimension bound used in
    \cite[Part~I]{TNLean2026OneDimensional}.
\end{theorem}
\begin{proof}
    \leanok
    \uses{lem:rank_one_entropy_zero}
    The reduced state on $A$ is positive semidefinite with trace one.
    Its space has dimension one, so its rank is at most one and its
    von Neumann entropy vanishes.
\end{proof}


\begin{theorem}[The one-dimensional area law on a chain]
    \label{thm:chain_area_law}
    \lean{
        TNLean.MPS.AreaLaw.chainIntervalEntropyBound,
        TNLean.MPS.AreaLaw.chainDiscountConst_nonneg,
        TNLean.MPS.AreaLaw.chainIntervalEntropyBound_nonneg,
        TNLean.MPS.AreaLaw.regionEntropy_interval_le,
        TNLean.MPS.AreaLaw.chainShellConst,
        TNLean.MPS.AreaLaw.chainShellCount_le,
        TNLean.MPS.AreaLaw.chainAreaLawConst,
        TNLean.MPS.AreaLaw.chainIntervalEntropyBound_le,
        TNLean.MPS.AreaLaw.regionEntropy_le_mul_card_boundaryEdges,
        TNLean.MPS.AreaLaw.chainAreaLaw
    }
    \leanok
    \uses{def:chain_model}
    Let $\Omega$ be a gapped ground vector, with gap $\Delta>0$, of a
    finite-range Hamiltonian with term norms at most $J\ge0$ on the chain of
    $L$ sites with local dimension $q\ge1$. Every interval $X$ satisfies
    \begin{align}
        S(X)\le C_{\mathrm{int}}(q,R,J,\Delta)
        =8mR\log q+\frac{4096\,e\,\vartheta B_R}{m}+m\log2
        \le C_0(q,R)\,(1+J/\Delta),
    \end{align}
    with
    $C_0(q,R)=\bigl(8R\log q+\log2\bigr)\bigl(5+32N_Rq^{2(R+1)}+128B_R\bigr)
        +1024\,e\,B_R$, and every subset $A$ of the chain satisfies
    $S(A)\le C_{\mathrm{int}}\,|\partial A|$. Hence the constant
    $C_0(q,R)(1+J/\Delta)$, chosen before $L$, the Hamiltonian, the ground
    vector and the region, bounds $S(A)/|\partial A|$ for every gapped chain.
    This is the connected case of Theorem~A; the general case is
    Theorem~\ref{thm:domain_area_law}.
\end{theorem}
\begin{proof}
    \leanok
    \uses{lem:chain_buffer_discount, lem:chain_entropy_inequalities}
    The discount of Lemma~\ref{lem:chain_buffer_discount} and the Araki--Lieb
    inequality of Lemma~\ref{lem:chain_entropy_inequalities} give
    $S(X)-2S(T)\le S(X\cup T)-S(T)\le\frac12S(X)+\frac12C_{\mathrm{disc}}$,
    hence $S(X)\le4S(T)+C_{\mathrm{disc}}$, and $S(T)\le|T|\log q\le
        2mR\log q$. The linear bound uses
    $m\le(5+32N_Rq^{2(R+1)}+128B_R)(1+\vartheta)$, which follows from
    $\sqrt{(1+\vartheta)B_R}\le(1+\vartheta)B_R$, and $m\ge4$. The subset
    bound is the second part of Lemma~\ref{lem:chain_entropy_inequalities}.
\end{proof}

\section{Finite sets of integers}

\begin{definition}[Finite domains, supports and admissible Hamiltonians]
    \label{def:domain_model}
    \lean{
        TNLean.MPS.AreaLaw.sitePos,
        TNLean.MPS.AreaLaw.sitePos_mem,
        TNLean.MPS.AreaLaw.exists_sitePos_eq,
        TNLean.MPS.AreaLaw.sub_le_sitePos_sub,
        TNLean.MPS.AreaLaw.val_eq_add_one_of_sitePos_eq_add_one,
        TNLean.MPS.AreaLaw.domainGraph,
        TNLean.MPS.AreaLaw.sitePos_sub_val_of_walk,
        TNLean.MPS.AreaLaw.IsDomainSupport,
        TNLean.MPS.AreaLaw.DomainSupport,
        TNLean.MPS.AreaLaw.IsDomainSupport.exists_walk,
        TNLean.MPS.AreaLaw.IsDomainSupport.isChainSupport,
        TNLean.MPS.AreaLaw.IsDomainSupport.subset_or_disjoint_of_gap,
        TNLean.MPS.AreaLaw.DomainHamiltonian,
        TNLean.MPS.AreaLaw.DomainHamiltonian.operator,
        TNLean.MPS.AreaLaw.DomainHamiltonian.toChain,
        TNLean.MPS.AreaLaw.DomainHamiltonian.toChain_operator,
        TNLean.MPS.AreaLaw.domainBoundaryEdges,
        TNLean.MPS.AreaLaw.domainRegion,
        TNLean.MPS.AreaLaw.IsDomainEdge,
        TNLean.MPS.AreaLaw.card_domainBoundaryEdges
    }
    \leanok
    \uses{def:chain_model}
    A finite set $\Lambda\subset\mathbb Z$ carries the induced path graph, whose
    edges join consecutive integers that both lie in $\Lambda$, and its graph
    distance $d_\Lambda$. Its sites are labelled $0,\ldots,|\Lambda|-1$ in
    increasing order, each carrying $\mathbb C^q$. A \emph{support of range $R$}
    in $\Lambda$ is a nonempty set of sites of $d_\Lambda$-diameter at most $R$.
    Along a path of the induced graph the position of a site minus its label is
    constant, so a support lies between two consecutive gaps of $\Lambda$ and
    is a support of range $R$ of the chain of $|\Lambda|$ sites. An
    \emph{admissible Hamiltonian} assigns to every support $X$ of range $R$ in
    $\Lambda$ a Hermitian term $h_X$ acting as the identity off $X$ with
    operator norm at most $J$, zero terms allowed; with zero terms on the
    remaining supports of the chain it is a finite-range Hamiltonian on the
    chain with the same sum. The \emph{edge boundary} $\partial_\Lambda A$ of a
    set $A\subseteq\mathbb Z$ is the set of edges $\{x,x+1\}\subseteq\Lambda$
    with exactly one endpoint in $A$; these edges correspond to the boundary
    edges, on the chain, of the sites in $A$ that join sites at consecutive
    integers.
\end{definition}

\begin{lemma}[A cut that splits no term]
    \label{lem:zero_entropy_cut}
    \lean{
        TNLean.MPS.AreaLaw.isSupportedOn_sum,
        TNLean.MPS.AreaLaw.regionEntropy_eq_zero_of_forall_subset_or_disjoint
    }
    \leanok
    Let $\Omega$ be a gapped ground vector, with gap $\Delta>0$, of a sum of
    Hermitian terms on a finite tensor product, each supported inside a set $B$
    of sites or inside its complement. Then $S_\Omega(B)=0$.
\end{lemma}
\begin{proof}
    \leanok
    \uses{thm:area_law_product_ground_state}
    In the coordinates of the cut $B\,|\,B^c$ the sum is
    $H_1\otimes1+1\otimes H_2$ with Hermitian $H_1$ and $H_2$. With least
    eigenvalues $a$, $b$ and unit eigenvectors $u$, $v$, the inequality
    $H\ge a+b$ gives $E_0\ge a+b$, and the gap inequality evaluated on
    $u\otimes v$ gives $\Delta(1-|\langle\Omega,u\otimes v\rangle|^2)\le
        a+b-E_0\le0$. Hence $\Omega$ is a multiple of $u\otimes v$, and its
    reduced state on $B$ is pure; this is
    Theorem~\ref{thm:area_law_product_ground_state}.
\end{proof}

\begin{lemma}[Flipping a region of zero entropy]
    \label{lem:zero_entropy_flip}
    \lean{
        TNLean.MPS.AreaLaw.regionEntropy_union_add_inter_le,
        TNLean.MPS.AreaLaw.regionEntropy_sdiff_eq_of_eq_zero,
        TNLean.MPS.AreaLaw.regionEntropy_union_eq_of_eq_zero,
        TNLean.MPS.AreaLaw.regionEntropy_symmDiff_eq_of_eq_zero,
        TNLean.MPS.AreaLaw.exists_mem_val_eq_iff,
        TNLean.MPS.AreaLaw.boundaryEdges_symmDiff_interval,
        TNLean.MPS.AreaLaw.exists_regionEntropy_eq_boundaryEdges_eq_symmDiff
    }
    \leanok
    \uses{lem:chain_entropy_inequalities}
    Let $\Omega$ be a unit vector on the chain of $L$ sites and $P$ a set of
    sites with $S(P)=0$. Then $S(A\triangle P)=S(A)$ for every set of sites
    $A$. If $P=[0,p+1)$ with $p+1<L$, the boundary edges of $A\triangle P$ are
    those of $A$ with the edge $\{p,p+1\}$ added or removed. Consequently, if
    $G$ is a set of edges $\{p,p+1\}$ whose prefixes $[0,p+1)$ all have zero
    entropy, every $A$ has a set of sites $A'$ with $S(A')=S(A)$ whose boundary
    edges are those of $A$ toggled exactly at the edges of $G$.
\end{lemma}
\begin{proof}
    \leanok
    \uses{lem:chain_entropy_inequalities}
    For $X\subseteq P$ the Araki--Lieb inequality gives
    $S(X)\le S(P)+S(P\setminus X)=S(P\setminus X)$ and symmetrically, so
    $S(X)=S(P\setminus X)$. For $X\subseteq P$ and $Y$ disjoint from $P$,
    strong subadditivity in the form
    $S(P\cup Y)+S(X)\le S(P)+S(X\cup Y)$ and the Araki--Lieb inequality
    $S(Y)\le S(P\cup Y)+S(P)$ give $S(X)+S(Y)\le S(X\cup Y)$, and
    subadditivity gives the reverse. Writing $A=(A\cap P)\cup(A\setminus P)$
    and $A\triangle P=(P\setminus(A\cap P))\cup(A\setminus P)$, both entropies
    equal $S(A\cap P)+S(A\setminus P)$. Flipping $[0,p+1)$ reverses membership
    at every site up to $p$, so it changes the status of the edge
    $\{p,p+1\}$ only. The last claim follows by induction on $G$.
\end{proof}

\begin{theorem}[Theorem A: the one-dimensional area law]
    \label{thm:domain_area_law}
    \lean{
        TNLean.MPS.AreaLaw.regionEntropy_le_mul_card_domainBoundaryEdges,
        TNLean.MPS.AreaLaw.domainAreaLaw
    }
    \leanok
    \uses{def:domain_model}
    Let $\Omega$ be a gapped ground vector, with gap $\Delta>0$, of an
    admissible Hamiltonian with term norms at most $J\ge0$ on a finite set
    $\Lambda\subset\mathbb Z$ with local dimension $q\ge1$. Every
    $A\subseteq\mathbb Z$ satisfies
    $S_\Omega(A\cap\Lambda)\le C_{\mathrm{int}}(q,R,J,\Delta)\,
        |\partial_\Lambda A|\le C_0(q,R)\,(1+J/\Delta)\,|\partial_\Lambda A|$,
    with the constants of Theorem~\ref{thm:chain_area_law}.
\end{theorem}
\begin{proof}
    \leanok
    \uses{thm:chain_area_law, lem:zero_entropy_cut, lem:zero_entropy_flip}
    The Hamiltonian is a finite-range Hamiltonian on the chain of $|\Lambda|$
    sites with the same sum, so Theorem~\ref{thm:chain_area_law} bounds the
    entropy of every set of sites by $C_{\mathrm{int}}$ times its number of
    boundary edges on the chain. At an edge $\{p,p+1\}$ of the chain whose
    sites are at nonconsecutive integers no support crosses, so the prefix
    $[0,p+1)$ has zero entropy by Lemma~\ref{lem:zero_entropy_cut}. Let $G$ be
    the set of such edges that are boundary edges of the sites in $A$. By
    Lemma~\ref{lem:zero_entropy_flip} there is a set of sites with the same
    entropy whose boundary edges on the chain are the remaining ones, which
    correspond to the edges of $\partial_\Lambda A$.
\end{proof}

\section{Matrix product states and sequential truncation}

\begin{definition}[Open-boundary matrix product states and cut matrices]
    \label{def:chain_open_mps}
    \lean{
        TNLean.MPS.AreaLaw.IsOpenMPS,
        TNLean.MPS.AreaLaw.IsOpenMPS.smul,
        TNLean.MPS.AreaLaw.leftBlock,
        TNLean.MPS.AreaLaw.mem_leftBlock,
        TNLean.MPS.AreaLaw.leftBlock_eq_interval,
        TNLean.MPS.AreaLaw.cutMatrix,
        TNLean.MPS.AreaLaw.cutMatrix_apply
    }
    \leanok
    \uses{def:chain_model}
    Let $L,q,D$ be natural numbers. A vector $\psi$ on $L$ sites of local
    dimension $q$ is an \emph{open-boundary matrix product state of bond
        dimension at most $D$} when there are natural numbers $r_0,\ldots,r_L$
    with $r_0=r_L=1$ and $r_k\leq D$ for $0<k<L$, a scalar $c\in\mathbb C$,
    and matrices $A_i(s)\in\mathbb C^{r_i\times r_{i+1}}$ such that
    \begin{align}
        \psi(x_0,\ldots,x_{L-1})
         & =c\bigl(A_0(x_0)\cdots A_{L-1}(x_{L-1})\bigr)_{0,0}. \notag
    \end{align}
    The empty product is the identity on $\mathbb C$, so at $L=0$ this
    condition reads $\psi(\varnothing)=c$. Only the internal bond dimensions
    are bounded by $D$. Scaling $\psi$ scales $c$.
    The \emph{left block} of $k$ sites is the interval $[0,k)$, and the
    \emph{cut matrix} at $k$ has the configurations of the left block as rows
    and those of the remaining sites as columns; its rank is the Schmidt
    rank across the cut.
\end{definition}

\begin{theorem}[Bounded internal cut ranks give a matrix product state]
    \label{thm:chain_mps_of_cut_rank}
    \lean{
        TNLean.MPS.AreaLaw.isOpenMPS_of_internal_cutMatrix_rank_le,
        TNLean.MPS.AreaLaw.isOpenMPS_of_cutMatrix_rank_le,
        Submodule.exists_fin_span_of_finrank_le,
        TNLean.MPS.AreaLaw.glueLeft,
        TNLean.MPS.AreaLaw.sliceMatrix,
        TNLean.MPS.AreaLaw.sliceMatrix_apply,
        TNLean.MPS.AreaLaw.rank_sliceMatrix_le,
        TNLean.MPS.AreaLaw.sliceSpace,
        TNLean.MPS.AreaLaw.finrank_sliceSpace_le,
        TNLean.MPS.AreaLaw.glueLeft_snoc,
        TNLean.MPS.AreaLaw.slice_mem_sliceSpace
    }
    \leanok
    Let $L,q,D$ be natural numbers and let $\psi$ be a vector on $L$ sites
    of local dimension $q$. If every internal cut matrix, at $0<k<L$, has
    rank at most $D$, then $\psi$ is an open-boundary matrix product state
    of bond dimension at most $D$. In particular, the conclusion holds
    when every cut matrix has rank at most $D$.
\end{theorem}
\begin{proof}
    \leanok
    \uses{def:chain_open_mps}
    Suppose first that $L\geq2$ and $D\geq1$. The two boundary cut matrices
    have rank at most $1\leq D$. Let $V_k$ be the column space of the cut
    matrix at $k$, and choose a spanning family $u_k$ of $D$ vectors.
    Slicing at a fixed value of site $k$ maps $V_{k+1}$ into $V_k$.
    Expressing these slices in the chosen spanning families gives square
    matrices $A_k(s)$ and boundary vectors $v_L,v_R$ with
    $\psi(x)=v_L^{\mathsf T}A_0(x_0)\cdots A_{L-1}(x_{L-1})v_R$.
    Absorb $v_L^{\mathsf T}$ into the first matrix and $v_R$ into the last;
    these are distinct matrices because $L\geq2$. Their rectangular
    dimensions have singleton endpoints and internal dimensions $D$.
    If $L\geq2$ and $D=0$, the cut matrix at $1$ vanishes, hence $\psi=0$.
    Set every internal bond dimension to zero and take zero site matrices
    and $c=0$. If $L=0$, take $c=\psi(\varnothing)$. If $L=1$, take
    $c=1$ and the $1\times1$ matrices $A_0(s)_{0,0}=\psi(s)$.
    The assertion about all cuts follows by restricting to internal cuts.
\end{proof}

\begin{lemma}[Sequential truncation]
    \label{lem:chain_sequential_truncation}
    \lean{
        Matrix.schmidtCoeffMatrix_kronecker_one_mulVec,
        Matrix.toEuclideanLin_one_apply,
        Matrix.toEuclideanLin_mul_apply,
        Entropy.localLift_sub,
        Entropy.exists_localLift_eq_localLift_of_subset,
        Entropy.norm_sq_toEuclideanLin_localLift_eq_trace,
        Entropy.norm_toEuclideanLin_localLift_le,
        Entropy.trace_regionState_eq_one,
        Matrix.PosSemidef.trace_complement_spectralProjectionGE_le_surprisalTail,
        Matrix.PosSemidef.rank_spectralProjectionGE_exp_le,
        TNLean.MPS.AreaLaw.cutMatrix_localLift,
        TNLean.MPS.AreaLaw.rank_cutMatrix_localLift_le,
        TNLean.MPS.AreaLaw.rank_cutMatrix_zero_le,
        TNLean.MPS.AreaLaw.rank_cutMatrix_le_one_of_le,
        TNLean.MPS.AreaLaw.LeftBlockMatrix,
        TNLean.MPS.AreaLaw.cutTruncation,
        TNLean.MPS.AreaLaw.sequentialTruncation,
        TNLean.MPS.AreaLaw.leftBlock_mono,
        TNLean.MPS.AreaLaw.norm_sub_sequentialTruncation_le,
        TNLean.MPS.AreaLaw.exists_sequentialTruncation_eq_localLift,
        TNLean.MPS.AreaLaw.rank_cutMatrix_sequentialTruncation_le,
        TNLean.MPS.AreaLaw.exists_sequentialTruncation
    }
    \leanok
    \uses{def:chain_open_mps}
    Let $\psi$ be a unit vector on $L$ sites and let $w$ be a real number. For
    each cut $k$,
    $1\le k\le L-1$, let $S_k$ be the entropy of the first $k$ sites, let
    $\Pi_k$ be the spectral projector of their reduced state onto the
    eigenvalues at least $e^{-(S_k+w)}$, an operator on the first $k$ sites,
    and let $P_k=\Pi_k\otimes1$ be its extension by the identity on the
    remaining sites. Then the regional projector $\Pi_k$ has rank at most
    $e^{S_k+w}$, the mass $\|(1-P_k)\psi\|^2$ discarded by $P_k$ is at most
    the surprisal tail of the reduced state beyond $w$, and the vector
    $\Phi=P_1P_2\cdots P_{L-1}\psi$ has cut rank at most
    $\max\{1,e^{S_k+w}\}$ across every cut $k$ of the chain, the trivial
    cuts at its ends included; it satisfies
    \begin{align}
        \|\psi-\Phi\|\le\sum_{k=1}^{L-1}\sqrt{\text{tail}_k(w)} .
    \end{align}
    This is the threshold form of the sequential truncation, Proposition~7.1
    of \cite{TNLean2026OneDimensional}: the note projects every cut onto its
    $D$ largest eigenvalues, or onto all of them when the cut has fewer than
    $D$, for an arbitrary bond dimension $D$, with the error
    $\sum_k\sqrt{\epsilon_k(D)}$; that form is not formalized, and
    Theorem~B uses only the threshold form; see
    \cite{gap:tnlean_one_dimensional_threshold_truncation}.
\end{lemma}
\begin{proof}
    \leanok
    \uses{def:chain_open_mps}
    Write $M_k(\psi)$ for the cut matrix of $\psi$ at $k$, with the
    configurations of the first $k$ sites as rows. For an operator $\Pi$ on
    the first $k$ sites, the extension $\Pi\otimes1$ acts on it by left
    multiplication,
    \begin{align}
        M_k\bigl((\Pi\otimes1)\psi\bigr)=\Pi\,M_k(\psi),
    \end{align}
    so the cut rank of $(\Pi\otimes1)\psi$ is at most the rank of $\Pi$; the
    rank of the extended operator itself, which is the rank of $\Pi$ times the
    dimension of the remaining sites, plays no role. The projector on the
    smallest block is
    applied last, and the product of the projectors on blocks of size at most
    $k$ is one extension of an operator on the first $k$ sites ending in
    $\Pi_k$. The distance follows by induction from
    $\psi-P_k\Phi'=(1-P_k)\psi+P_k(\psi-\Phi')$ and the contraction property
    of a projector, and $\|(1-P_k)\psi\|^2$ is the trace of $1-\Pi_k$ against
    the reduced state, the mass of the eigenvalues below the threshold, which
    is at most the surprisal tail.
\end{proof}

\begin{theorem}[Theorem B: polynomial matrix product state approximation]
    \label{thm:chain_mps_approximation}
    \lean{
        TNLean.MPS.AreaLaw.tailPrefactorRoot,
        TNLean.MPS.AreaLaw.tailPrefactorRoot_pos,
        TNLean.MPS.AreaLaw.tailPrefactorRoot_sq,
        TNLean.MPS.AreaLaw.one_lt_tailPrefactorRoot,
        TNLean.MPS.AreaLaw.chainTruncationWidth,
        TNLean.MPS.AreaLaw.chainTruncationWidth_nonneg,
        TNLean.MPS.AreaLaw.exp_neg_chainTruncationWidth,
        TNLean.MPS.AreaLaw.chainBondDimension,
        TNLean.MPS.AreaLaw.exists_isOpenMPS_norm_sub_le,
        TNLean.MPS.AreaLaw.chainBondDimension_le_mul_rpow
    }
    \leanok
    \uses{def:chain_model, def:chain_open_mps}
    Let $\Omega$ be a gapped ground vector, with gap $\Delta>0$, of a
    finite-range Hamiltonian with term norms at most $J\ge0$ on the chain of
    $L\ge1$ sites with local dimension $q\ge1$. Put
    $s=\sqrt{(1+J/\Delta)B_R}$, $c_0=\sqrt{2e^{e/2}}$ and
    \begin{align}
        w(L)=64\,s\log\bigl(2L^2c_0\bigr),\qquad
        D(L)=\bigl\lfloor\exp\bigl(C_{\mathrm{int}}+w(L)\bigr)\bigr\rfloor
        \le e^{C_{\mathrm{int}}}(2c_0)^{64s}\,L^{128s}.
    \end{align}
    There is a normalized open-boundary matrix product state $\Phi$ of bond
    dimension at most $D(L)$ with $\|\Phi-\Omega\|\le1/L$.
\end{theorem}
\begin{proof}
    \leanok
    \uses{thm:chain_area_law, lem:chain_marginal_tail,
        lem:chain_sequential_truncation, thm:chain_mps_of_cut_rank,
        lem:chain_crossing_count}
    By Theorem~\ref{thm:chain_area_law} the entropy of every left block, an
    interval $[0,k)$, is at most $C_{\mathrm{int}}$, and by the tail bound of
    Lemma~\ref{lem:chain_marginal_tail} with the budget $B_R$ of
    Lemma~\ref{lem:chain_crossing_count} its surprisal tail beyond $w$ is at
    most $2e^{e/2}e^{-w/(32s)}$. The sequential truncation of
    Lemma~\ref{lem:chain_sequential_truncation} at width $w(L)$ gives cut
    ranks at most $e^{C_{\mathrm{int}}+w(L)}$ and distance at most
    $(L-1)\cdot\frac1{2L^2}\le\frac1{2L}$. Normalizing moves the vector by at
    most the same amount, so the distance is at most $1/L$, and
    Theorem~\ref{thm:chain_mps_of_cut_rank} exhibits the matrix product
    structure, which survives the normalization.
\end{proof}
